% Part: incompleteness
% Chapter: theories-computability
% Section: computably-axiomatizable

\documentclass[../../../include/open-logic-section]{subfiles}

\begin{document}

\olfileid{inc}{tcp}{cax}

\olsection{\printtoken{S}{axiomatizable} তত্ত্ব}

কোনো তত্ত্ব~$\Th{T}$-এর স্বতঃসিদ্ধসমষ্টি~$A$ গণনসাধ্য হলে তত্ত্বটিকে
\emph{!!{axiomatizable}} বলা হয়। ($A$-কে $\Th{T}$-এর স্বতঃসিদ্ধসমষ্টি
বলার অর্থ $T = \Setabs{!A}{A \Proves !A}$।) স্বাভাবিক সংখ্যার যেকোনো
``যুক্তিসঙ্গত'' স্বতঃসিদ্ধায়নের এই ধর্ম থাকবে। বিশেষত, সসীম
স্বতঃসিদ্ধসমষ্টিবিশিষ্ট প্রতিটি তত্ত্ব !!{axiomatizable}।

\begin{lem}
ধরি $\Th{T}$ !!{axiomatizable}। তাহলে $\Th{T}$
!!{computably
enumerable}।
\end{lem}

\begin{proof}
ধরি $A$ হলো $\Th{T}$-এর একটি গণনসাধ্য স্বতঃসিদ্ধসমষ্টি। $!A \in T$
হলে তা শনাক্ত করতে স্বতঃসিদ্ধগুলি থেকে $!A$-এর কোনো !!a{derivation}
খুঁজতে থাকলেই হয়।
% BN-SRC-238 TARGET-CORRECTION-BEGIN
\par\smallskip\noindent\textnormal{\small\textbf{উৎস-স্পষ্টীকরণ (BN-SRC-238):}
হিমায়িত পাঠে এই একমুখী অনুসন্ধানকে ``$!A \in T$ কি না নির্ধারণ'' বলা
হয়েছে। সদস্য না হলে প্রমাণ-অনুসন্ধান থামার নিশ্চয়তা নেই; তাই লেমার
c.e. সিদ্ধান্তের সঙ্গে মিলিয়ে বাংলা পাঠে কেবল সদস্য হলে শনাক্ত করার কথা
বলা হয়েছে।}
% BN-SRC-238 TARGET-CORRECTION-END

কথাটি আরেকভাবে বলা যায়: $!A$ ঠিক তখনই $\Th{T}$-তে থাকে যখন $A$-তে
$!B_1$,~\dots, $!B_k$ স্বতঃসিদ্ধের এমন একটি সসীম তালিকা এবং প্রথম-ক্রমের
যুক্তিবিদ্যায় $(!B_1 \land \dots \land !B_k) \lif !A$-এর কোনো
!!a{derivation} থাকে। কিন্তু আমরা ইতিমধ্যে জানি, ``এমন কিছু আছে যাতে
\dots'' আকারের সংজ্ঞাবিশিষ্ট যেকোনো সেট !!{c.e.}, যদি ``\dots''-এর
পরের শর্তটি গণনসাধ্য হয়।
\end{proof}

\end{document}
